\section{Related Work}
\label{sec:related}

\paragraph{Correspondence-based 6D pose estimation.}
Most leading RGB methods on BOP predict 2D--3D correspondences and solve for the pose with
PnP. EPOS~\cite{hodan2020epos} predicts dense surface fragments, CDPN~\cite{li2019cdpn}
disentangles rotation (from correspondences) and translation (regressed),
ZebraPose~\cite{su2022zebrapose} encodes surface points with hierarchical binary codes, and
GDR-Net / GDRNPP~\cite{wang2021gdrnet,liu2022gdrnpp} learn the PnP step from geometric
intermediate maps. These methods top the PBR-only RGB leaderboards, but most of them train one
network per object, and their inference includes PnP-RANSAC or several stages. In contrast, we
regress the pose directly with one network shared by all objects.

\paragraph{Render-and-compare and refinement.}
CosyPose~\cite{labbe2020cosypose}, SurfEmb~\cite{haugaard2022surfemb} and
MegaPose~\cite{labbe2022megapose} refine an initial estimate by rendering the object model and
comparing it with the image, iteratively. Refinement improves accuracy, but it multiplies the
inference time (up to seconds per image) and requires a renderer at test time. We use no
refinement; our estimates could serve as initialization for such a stage.

\paragraph{Direct regression and Transformers.}
Direct methods regress the rotation and translation from image features in a single forward
pass. The choice of the targets matters: continuous rotation representations~\cite{zhou2019continuity},
allocentric (apparent) rotations~\cite{kundu20183drcnn}, and scale-invariant translation
targets~\cite{li2019cdpn,wang2021gdrnet}. PViT-6D~\cite{stapf2023pvit6d} regresses the pose
of each detection with a multi-scale ViT (MViTv2) shared by all objects: separate rotation and
translation tokens, initialized with a learned embedding of the detected class, feed two MLP
heads; it introduces a size-invariant depth target that normalizes the depth by the object
diameter, which we also use (\cref{sec:method:translation}), and reaches 0.73--0.74 AR on LM-O
with PBR training. PoseDETR and RACE-6D (BOP leaderboard submissions) predict instances and poses
in a single stage at about 12\,ms per image. FastPose-ViT~\cite{ancey2026fastposevit} regresses normalized
translation targets and the apparent rotation from object crops with a ViT and a single linear
head, and runs in real time on a Jetson Orin Nano, but for one object and one camera only. We
extend it to many objects and cameras with a plain ViT and a single linear layer, and evaluate it
on BOP. \todo{Sharpen the difference with PViT-6D (plain ViT-B and one linear layer vs MViTv2 with
pose tokens and MLP heads; symmetric loss with offsets; deployment study) and cite PoseDETR /
RACE-6D papers when available.}

\paragraph{Fast object detectors.}
Two-stage pose pipelines depend on a 2D detector. The BOP default detections for the 2023
challenge come from the GDRNPP detector~\cite{liu2022gdrnpp}, which takes about 80\,ms per image.
Real-time DETR variants such as LW-DETR~\cite{chen2024lwdetr} make the detector much cheaper; we
train a multi-class LW-DETR on the same PBR images to obtain a fully real-time pipeline.
